\section{Discussion of the Research Gap:} What \emph{Could} We Design for When We Design Smart Buildings
\label{sec-ch2:discussion}
Our review identified 11 unique experiences and 20 distinct design mechanisms, reflecting the broad, experimental, and evolving nature of smart building research. While this diversity highlights a field still in its formative stages, the scattered approaches indicate an opportunity for consolidation and theory-building to provide clearer guidance. In the following, we point out research areas that we consider particularly relevant and promising for further development in the domain of smart building design. 

\paragraph{More-Than-Human Perspectives in Smart Buildings} Our review did not find more-than-human perspectives in smart buildings. These perspectives challenge the traditional user-centred design by emphasising environments where human and non-human actors are examined collectively~\cite{giaccardi2020technology, coskun2022more, wakkary2021things}. Within smart buildings, this could mean considering their environmental impact during human inhabitation as well as times of non-occupancy (e.g. during off-peak hours or after abandonment). A more-than-human perspective could guide technology to sustain relationships with non-human actors during these phases, by developing safety measures for local birds that may crash into building fa\c{c}ades~\cite{boone2024embodied}. Furthermore, smart buildings could engage with local ecosystems by adopting \textit{Biophilic Design} patterns, such as incorporating wildlife corridors to allow animal movement and reconnect fragmented habitats in urban areas~\cite{zellmer2022urban}. \textit{Empathic Design} mechanisms could foster a sense of ecological responsibility through the integration of ``animalistically behaving'' robotic components to enhance connections with local wildlife~\cite{gil2015biophilic}. Smart building devices could actively ``negotiate'' climate levels in office rooms with occupants, prioritising environmental considerations~\cite{loh2010please}.

Critically, more-than-human perspectives also raise ethical and practical questions. Who decides the ``priorities'' of non-human actors (such as the living bio-diversity in building ecosystems) in a building? How do we ensure that smart systems designed for ecological sensitivity do not inadvertently impose anthropocentric values onto non-human entities?  These considerations highlight the need for interdisciplinary collaboration to integrate ecological, ethical, and technological insights into smart building design. More-than-human co-design has already become a part of smart city discourses, considering relationships between humans and non-human actors including plants, animals, insects, as well as soil, water, and other ecological features~\cite{heitlinger2024designing, fieuw2022towards}. Drawing inspiration from these strategies, co-design approaches could provide an opportunity to engage ecological stakeholders (e.g. conservationists, environmental scientists) alongside participants~\cite{heitlinger2024designing}. This could ensure that more-than-human perspectives are meaningfully integrated into the design and evaluation of smart buildings.



\paragraph{Comfort as a Bodily Experience}
While somaesthetic approaches were observed in addressing \textbf{Emotional} and \textbf{Bodily Experiences}, \textbf{Comfort} continues to be predominantly defined by traditional metrics such as thermal and visual conditions, overlooking the embodied dimensions of how we experience the built environments we inhabit. The absence of \textit{Somaesthetic Design} in \textbf{Comfort} studies indicates an opportunity to rethink comfort through the lens of ``design with the body''~\cite{hook2018designing}. This could involve designing smart buildings that account for not just environmental parameters, but also the ways in which bodily postures, movements, and somatic experiences contribute to a holistic well-being~\cite{chappells2005debating}. Incorporating somaesthetics into comfort raises questions about what it means to ``live with our bodies'' in smart buildings. This approach may even question the very foundations of comfort studies. For example, should building design create opportunities for (mild) discomfort---as a human natural state---to promote mindfulness through enhanced bodily awareness~\cite{tschaepe2021somaesthetics}? This would be in contrast to projects like SonicTaiChi which integrates pleasing sound and slow movement to promote restoration~\cite{jakovich2007particletecture}. Instead, could we further explore smart architectural elements to intentionally induce acoustic discomfort to provoke the recognition of stress, not to alleviate stress temporarily, but to provoke awareness of its underlying causes, encouraging deeper self-reflection and the pursuit of long-term solutions?~\cite{tschaepe2021somaesthetics}.




\paragraph{Social Justice and Activism} 
Beyond civic values, \textbf{Democratic Experiences} could be extended to include social justice themes such as feminism~\cite{mitchell2020human}, racial justice ~\cite{corbett2019engaging}, and activist design perspectives (such as anti-oppressive design~\cite{fox2016exploring}), that were notably limited. These omissions are striking, given ongoing discussions in HCI and technology design ~\cite{chordia2024social}. Feminist and activist lenses could be crucial for understanding how smart buildings can either reinforce or challenge power structures and inequalities~\cite{mitchell2020human}. For example, a feminist lens could question how personalisation features in smart buildings (e.g. temperature in shared office spaces) disproportionately impact gender relations~\cite{boys10there}. Additionally, spatial design research  highlights that \textbf{Sense of Orientation} differs across genders, i.e. a tendency to rely on mental maps and cardinal directions, as opposed to favouring route-based navigation using landmarks and sequential cues~\cite{munion2019gender}. A feminist perspective could investigate how smart building features---such as adaptive wayfinding systems or customisable navigation aids---could support diverse cognitive strategies~\cite{kutlu2005sense, coluccia2004gender}.

Neurodiversity is another under-explored area in smart building design. Technologies could better accommodate sensory differences in individuals by incorporating features like adjustable lighting and soundscapes. For example, personalised smart lighting can benefit individuals with ADHD to retain attention and cognitive performance~\cite{bagheri2024visual}. Similarly, adapting the design of wayfinding systems can lead individuals with autism on routes that avoid spaces with certain disturbing environmental features (e.g. sound or odours) \cite{black2022considerations}. Despite growing recognition, the effective implementation of neurodiverse considerations in smart building technologies remains an open challenge~\cite{bagheri2024visual, bruyere2022neurodiversity, kwon2023interior}.

In the context of cultural activism, drawn from urban studies~\cite{buser2013cultural}, smart buildings could serve as platforms for creative practices aimed at social and political change. Interactive art installations or adaptive architectural features could spotlight social injustices such as unequal access to housing or energy resources~\cite{loszchuk2020combining}. Smart buildings could also act as mediums for cultural activism by creating spaces that foster collective storytelling, resistance, or solidarity among communities. Through mechanisms such as \textit{Playful Design} or \textit{Community and Social Connections}~\cite{shepard2011sentient}, smart building design could support activist practices like hosting protests or interactive public events that raise awareness of societal inequities. Applying these perspectives could re-direct smart building research towards new design horizons shaped by principles of social justice~\cite{dombrowski2016social}. 


\paragraph{Evaluating the Smart Building Experience} 
Through the process of this scoping review, we observed a limited emphasis of papers on the evaluation of smart building experiences. This raises the question of whether new assessment frameworks, metrics, and methodologies are needed to be developed. Traditional metrics like objective comfort measurements do not fully capture the complexities of smart building interactions, especially as these environments become more immersive, interactive, and personalised. For instance, in the context of educational spaces, \textbf{Learning and Teaching Experiences} highlight the potential to create smart buildings that encourage collaboration, exploration, and knowledge-sharing~\cite{lui2014supporting}. However, evaluation of these experiences often relies on conventional measures, such as students' concentration~\cite{korozi2019shaping}. This may overlook the impact of learning spaces on more complex experiences such as emotion~\cite{gao2020n, gao2022understanding} and playfulness~\cite{ccetin2024learning} that are recognised to affect learning. Therefore, additional metrics and new methods of evaluation are required that can account for the subjectivity and complexity of human interactive experiences in built environments. 


At the same time, some of the targeted human experiences---such as \textbf{Bodily Experiences} or \textbf{Aesthetic Experiences}--- require developing qualitative methods to evaluate technological interventions that are grounded in the knowledge of human psychological and physiological responses to environments (e.g. environmental psychology~\cite{bower2019impact}). Currently, these experiences are captured in terms of descriptive qualities, such as the emotional resonance of a space, the sensory richness of an environment, or the harmony between occupants and their surroundings~\cite{schuricht2007freequent, diniz2008body}. While these qualitative methods can provide valuable insights into the lived and embodied aspects of smart buildings, the opportunity to combine them with physiological sensing and associated quantitative methods (e.g. ~\cite{moran2013using}) has been largely overlooked. As a next step, future research should systematically review and map existing evaluation methods for assessing human experiences in smart buildings to provide a more accurate picture of existing methods and determine the need for new frameworks and approaches. 

\paragraph{Anticipated Technological Integration in Smart Buildings} Current smart building technologies are primarily driven by sensors and actuators, with emerging technologies such as AI-enabled systems and extended reality yet to find their widespread practical application in built environments~\cite{moreno2024generative}. The incorporation of such technologies has however been discussed with caution and awareness around unintended consequences related to privacy, data ownership, and the risk of over-reliance on artificial environments at the expense of natural, sensory-rich experiences~\cite{moreno2024generative}. For example, while generative AI (genAI) offers possibilities for dynamic and adaptive environments, its substantial carbon footprint raises critical questions about sustainability~\cite{chien2023reducing}. With the building industry already contributing significantly to global emissions, combining these two areas requires careful consideration. While genAI may not be strictly necessary for many purposes, it is likely to be used by researchers and practitioners interested in experimenting with it. One promising avenue lies in leveraging more lightweight genAI models~\cite{mei2023lightlm} (e.g. for \textbf{Aesthetic Experiences}, which could enhance sensory engagement and atmosphere) without having to hinge on data-intensive AI processes. By leveraging ``scalable'' environmental elements like light, sound, and spatial configurations~\cite{deng2019diffusive, ishiguro2014ubisonus}, lightweight genAI could support the creation of dynamic and enriching environments with a reduced ecological footprint. Adaptive lighting and soundscapes powered by such models could offer evolving atmospheres tailored to user preferences while minimising energy consumption~\cite{zhao2017mediated}. Integrating \textbf{Aesthetic Experiences} with XR and lightweight genAI models could further enable novel forms of creative expression in smart buildings. These systems might craft immersive environments ``prompted'' by occupants based on their preferences to enable aspects such as creativity over purely functional or utilitarian goals in smart spaces \cite{ccetin2024learning}.

\paragraph{Anticipated Technological Integration in Smart Buildings} Current smart building technologies are primarily driven by sensors and actuators, with emerging technologies such as AI-enabled systems and extended reality yet to find a widespread practical application in built environments~\cite{moreno2024generative}. The incorporation of such technologies has however been discussed with caution and awareness around unintended consequences related to privacy, sustainability, and the risk of over-reliance on artificial environments at the expense of natural, sensory-rich experiences~\cite{moreno2024generative}. For example, while generative AI (genAI) offers possibilities for dynamic and adaptive environments, its substantial carbon footprint raises critical questions about sustainability~\cite{chien2023reducing}. With the building industry already contributing significantly to global emissions, combining these two areas requires careful consideration. While genAI may not be strictly necessary for many purposes, it is likely to be used by researchers and practitioners interested in experimenting with it. One promising avenue lies in leveraging more lightweight genAI models~\cite{mei2023lightlm} for \textbf{Aesthetic Experiences}, which could enhance sensory engagement and atmosphere without having to hinge on data-intensive AI processes. By leveraging ``scalable'' environmental elements like light, sound, and spatial configurations~\cite{deng2019diffusive, ishiguro2014ubisonus}, lightweight genAI could support the creation of dynamic and enriching environments with a reduced ecological footprint. Adaptive lighting and soundscapes powered by such models could offer evolving atmospheres tailored to user preferences while minimising energy consumption~\cite{zhao2017mediated}. Integrating \textbf{Aesthetic Experiences} with XR and lightweight genAI models could further enable novel forms of creative expression in smart buildings. These systems might craft immersive environments ``prompted'' by occupants based on their preferences to enable aspects such as creativity over purely functional or utilitarian goals in smart spaces~\cite{ccetin2024learning}. 


\section{Discussion of Limitations: Researcher Lens in Creation of Targeted Experiences and Design Mechanisms}
We acknowledge that our work aligns with the third paradigm in HCI, making the research process interpretive and shaped by our individual and collective perspectives \cite{harrison2011making}. Rather than offering definitive or universal insights, we present these findings as part of an ongoing discourse, encouraging reflection from the community. In developing the targeted experiences and design mechanisms, we acknowledge the influence of the researchers' lens in shaping how we identified commonalities and differences between the works. For instance, \textit{Somaesthetic Design} and \textbf{Bodily Experiences} reflect the ongoing interests within the HCI community, while \textit{Biophilic Design} might be influenced by the authors' interests in multi-sensory connections to nature in smart buildings. The creation of the \textbf{Explorative and/or Future and Speculative} experience also draws from our background in architecture research and IxD, capturing how digital technologies transform physical spaces. Therefore, we acknowledge that our positionality as researchers in the fields of HCI and architecture influenced how we categorised and interpreted the material.

Some corpus papers considered what could be interpreted as experiences of care in the context of health, care-giving, and spaces that meet users' needs (e.g. \cite{hosseini2023integrating, sowles2021introducing}). We did not establish a distinct experience around this concept, recognising that these aspects are inherently embedded within the other experiences we identified. For instance, many works within \textbf{Emotional Experiences} and \textbf{Comfort} are predominantly underpinned by health and care-giving considerations~\cite{alavi2017comfort, zhong2020hilo}. Mechanisms like \textit{Community and Social Connection} and \textit{Technological Immersion} also indirectly foster care by promoting socialisation and immersive health-related interactions. As care-giving and health permeate multiple smart building experiences, future research may explore these topics more explicitly using our dataset. Given this, we believe that care-giving and health are woven into multiple dimensions of smart building experiences. However, researchers focused on these topics may choose to view our dataset through a different lens to explore care-giving more explicitly.

Our particular search keywords may also have impacted our overall set of works. To narrow our search, we included the term ``interac*'' to focus on human interaction. This may have inadvertently excluded relevant studies from outside the HCI community that do not explicitly mention interaction, yet still address human experiences in smart buildings, although we found no direct evidence that this was the case. Finally, we categorised papers based on primary and overarching experiences that they targeted. Future research expanding our work to capture all experiences, i.e., including secondary and tangential ones, may uncover further targeted human experiences. Rather, this categorisation enabled us to elicit distinct features of the smart building research landscape and to conduct a deeper exploration of specific core experiences and accompanying design mechanisms. 
